\makesectionpage[figures/quiz.png]{Datenaufbereitung}

\begin{frame}[fragile]{Arten von Daten: Kategoriale Daten}
\begin{columns}
\column{0.47\textwidth}
\begin{itemize}
\item Datensätze enthalten \textbf{nicht nur numerische Werte}
\item Beispiel: Spalte \texttt{cut}
\item Beschreibt Qualität des Schliffs
\item \textbf{Ordinale Kategorie (Rangfolge)}:
\begin{itemize}
\item Fair
\item Good
\item Very Good
\item Premium
\item Ideal
\end{itemize}
\item Wichtig für Datenanalyse und ML
\end{itemize}

\column{0.48\textwidth}
\begin{minted}{python}
import pandas as pd

df = pd.read_csv("diamonds.csv")
df.head()
\end{minted}

\vspace{0.2cm}

\small
\begin{tabular}{r l l r}
\textbf{carat} & \textbf{cut} & \textbf{color} & \textbf{price} \\
\hline
0.23 & Ideal & E & 326 \\
0.21 & Premium & E & 326 \\
0.23 & Good & E & 327 \\
0.29 & Premium & I & 334 \\
0.31 & Good & J & 335 \\
\end{tabular}
\end{columns}

\source{\cite{seemann2024mlaz}, DiamantenOrdinaleDaten.ipynb}
\end{frame}

\begin{frame}[fragile]{Arten von Daten: Nominale und unstrukturierte Daten}
\begin{columns}
\column{0.47\textwidth}
\begin{itemize}
\item Datensatz einer Kleinanzeigenplattform
\item Ziel: Prognose des Verkaufspreises
\item Spalte \texttt{fuelType}:
\begin{itemize}
\item Werte: benzin, diesel, hybrid, lpg
\item Keine natürliche Ordnung
\item Nominalskalierte Kategorie
\end{itemize}
\item Spalte \texttt{name}:
\begin{itemize}
\item Freitext (unstrukturierte Daten)
\item Enthält Zusatzinformationen
\item Vor Verarbeitung aufzubereiten
\end{itemize}
\end{itemize}

\column{0.48\textwidth}
\begin{minted}[fontsize=\footnotesize]{python}
import pandas as pd

df = pd.read_csv("./fahrzeug_preis.csv")
df.head()
\end{minted}

\vspace{0.2cm}

\footnotesize
\begin{tabular}{l l l l}
\textbf{price} & \textbf{year} & \textbf{fuelType} & \textbf{name} \\
\hline
1450 & 1997 & benzin & Toyota\_Toy \\
13100 & 2005 & benzin & R32\_tausch \\
4500 & 2008 & benzin & Toyota\_Yar \\
6000 & 2009 & diesel & 320\_Alpinw \\
3990 & 1999 & benzin & BMW\_318i\_E \\
\end{tabular}
\end{columns}

\source{\cite{seemann2024mlaz}, OneHotFahrzeugdaten.ipynb}
\end{frame}


\begin{frame}{Arten von Daten}
\begin{columns}
\column{0.35\textwidth}
Insgesamt haben wir die nebenstehenden Arten von Daten.

\column{0.60\textwidth}
\centering
\resizebox{\linewidth}{!}{%
\begin{tikzpicture}[
  every node/.style={font=\scriptsize},
  main/.style={draw, rounded corners=4pt, very thick, align=center,
               inner sep=7pt, text width=2.4cm,
               fill=PrimaryColor!12, font=\scriptsize\bfseries},
  sub/.style={draw, rounded corners=4pt, thick, align=left,
              inner sep=7pt, text width=4.2cm, fill=SecondaryColor!8},
  line/.style={thick, SecondaryColor!75!black, ->, >=stealth},
  trunk/.style={thick, SecondaryColor!75!black},
]

% --- Hauptknoten ---
\node[main, draw=PrimaryColor] (strukt)   at (0, 0.4)  {Strukturierte\\ Daten};
\node[main, draw=PrimaryColor] (unstrukt) at (0,-3.5)  {Unstrukturierte\\ Daten};

% --- Unterkategorien (strukturiert) ---
\node[sub, draw=SecondaryColor] (metrisch) at (4.8, 2.1) {%
  \textbf{Metrische Daten}\\[2pt]
  Zahlenwerte (Komma-Werte)\\[1pt]
  {\itshape\color{SecondaryColor!55!black} z.\,B. Abmessungen eines Diamanten}};

\node[sub, draw=SecondaryColor] (ordinal) at (4.8, 0.4) {%
  \textbf{Ordinale Daten}\\[2pt]
  Kategorien \emph{mit} Rangordnung\\[1pt]
  {\itshape\color{SecondaryColor!55!black} z.\,B. Reinheit, Farbe, Schliff}};

\node[sub, draw=SecondaryColor] (nominal) at (4.8,-1.3) {%
  \textbf{Nominale Daten}\\[2pt]
  Kategorien \emph{ohne} Ordnung\\[1pt]
  {\itshape\color{SecondaryColor!55!black} z.\,B. Kraftstofftyp}};

% --- Detail (unstrukturiert) ---
\node[sub, draw=SecondaryColor] (freitext) at (4.8,-3.5) {%
  \textbf{Freitext}\\[2pt]
  Vorverarbeitung nötig\\[1pt]
  {\itshape\color{SecondaryColor!55!black} z.\,B. Überschrift der Fahrzeuganzeige}};

% --- Verbindungen: verzweigter Konnektor ---
\draw[trunk] (strukt.east) -- (2.1,0.4);
\draw[trunk] (2.1, 2.1) -- (2.1,-1.3);
\draw[line]  (2.1, 2.1) -- (metrisch.west);
\draw[line]  (2.1, 0.4) -- (ordinal.west);
\draw[line]  (2.1,-1.3) -- (nominal.west);

\draw[line]  (unstrukt.east) -- (freitext.west);

\end{tikzpicture}}

\end{columns}
\source{\cite{seemann2024mlaz}}
\end{frame}


\begin{frame}{Nominale Daten: One-Hot Encoding}
\begin{columns}
\column{0.47\textwidth}
\begin{itemize}
  \item Viele ML-Algorithmen brauchen numerische Werte und können nicht mit nominalen Daten umgehen.
  \item Naive Nummerierung
        (\texttt{benzin}=0, \texttt{diesel}=1, \texttt{hybrid}=2, \texttt{lpg}=3)
        suggeriert \textbf{messbare Abstände}.
  \item Das würde unserem Algorithmus aber vermitteln, dass ein Benzin-Antrieb einem Diesel näher steht  als einem Hybrid. 
  
\end{itemize}

\column{0.48\textwidth}
\begin{itemize}
  \item \textbf{One-Hot Encoding:} je möglichem Wert eine eigene Spalte
        \begin{itemize}
          \item Wert \texttt{1}, wenn zutreffend
          \item sonst Wert \texttt{0}
        \end{itemize}
\end{itemize}

\medskip
\centering
\renewcommand{\arraystretch}{1.5}
\resizebox{\linewidth}{!}{%
\begin{tabular}{l cccc}
\rowcolor{PrimaryColor}
\textcolor{white}{\textbf{fuelType}} &
\textcolor{white}{\textbf{benzin?}} &
\textcolor{white}{\textbf{diesel?}} &
\textcolor{white}{\textbf{hybrid?}} &
\textcolor{white}{\textbf{lpg?}}\\
\cellcolor{PrimaryColor}\textcolor{white}{\textbf{benzin}} & 1 & 0 & 0 & 0\\
\rowcolor{SecondaryColor!12}
\cellcolor{PrimaryColor}\textcolor{white}{\textbf{diesel}} & 0 & 1 & 0 & 0\\
\cellcolor{PrimaryColor}\textcolor{white}{\textbf{hybrid}} & 0 & 0 & 1 & 0\\
\rowcolor{SecondaryColor!12}
\cellcolor{PrimaryColor}\textcolor{white}{\textbf{lpg}}    & 0 & 0 & 0 & 1\\
\end{tabular}}

\vspace{0.6em}
{\footnotesize\itshape Aus einer Spalte werden vier 0/1-Spalten.}
\end{columns}
\source{\cite{seemann2024mlaz}}
\end{frame}

\begin{frame}{Nominale Daten: One-Hot Encoding}
\begin{columns}
\column{0.47\textwidth}
\begin{itemize}
  \item Die \textbf{letzte Spalte} ist immer redundant.
  \item Ihr Wert lässt sich stets aus den übrigen Spalten ableiten
        ({\itshape alle 0 $\Rightarrow$ es bleibt nur \texttt{lpg}}).
  \item Redundanz stört bei der Modellbildung.
  \item[$\Rightarrow$] Diese Spalte wird \textbf{gelöscht}.
\end{itemize}

\column{0.48\textwidth}
\centering
\renewcommand{\arraystretch}{1.5}
\begin{tikzpicture}[remember picture]
\node[inner sep=0pt] (tab) {%
\footnotesize
\begin{tabular}{l cccc}
\rowcolor{PrimaryColor}
\textcolor{white}{\textbf{fuelType}} &
\textcolor{white}{\textbf{benzin?}} &
\textcolor{white}{\textbf{diesel?}} &
\textcolor{white}{\textbf{hybrid?}} &
\tikzmark{xtl}\textcolor{white}{\textbf{lpg?}}\\
\cellcolor{PrimaryColor}\textcolor{white}{\textbf{benzin}} & 1 & 0 & 0 & 0\\
\rowcolor{SecondaryColor!12}
\cellcolor{PrimaryColor}\textcolor{white}{\textbf{diesel}} & 0 & 1 & 0 & 0\\
\cellcolor{PrimaryColor}\textcolor{white}{\textbf{hybrid}} & 0 & 0 & 1 & 0\\
\rowcolor{SecondaryColor!12}
\cellcolor{PrimaryColor}\textcolor{white}{\textbf{lpg}}    & 0 & 0 & 0 & 1\tikzmark{xbr}\\
\end{tabular}};
\end{tikzpicture}

% --- rotes X ueber die letzte Spalte ---
\begin{tikzpicture}[overlay, remember picture]
  \coordinate (TL) at ($(pic cs:xtl)+(-0.25,0.45)$);
  \coordinate (BR) at ($(pic cs:xbr)+(0.30,-0.20)$);
  \draw[red, line width=2.2pt, opacity=0.85] (TL) -- (BR);
  \draw[red, line width=2.2pt, opacity=0.85] (TL -| BR) -- (BR -| TL);
\end{tikzpicture}

\vspace{0.4em}
{\footnotesize\itshape Spalte \texttt{lpg?} kann entfallen.}
\end{columns}
\source{\cite{seemann2024mlaz}}
\end{frame}


\begin{frame}[fragile]{Nominale Daten: One-Hot Encoding in Python}
\begin{columns}
\column{0.60\textwidth}
\begin{minted}[fontsize=\footnotesize]{python}
df = pd.get_dummies(df, columns=["fuelType"])
df = df.drop(labels=["fuelType_lpg"], axis=1)

df = df.drop(labels=["model", "name"], axis=1)
df.head()
\end{minted}

\column{0.35\textwidth}
\footnotesize
\textbf{Tipp:} One-Hot Encoding erzeugt schnell \textbf{viele Spalten}.
Nicht benötigte Spalten direkt mit \texttt{drop} entfernen. Das hält
den Datensatz übersichtlich.
\end{columns}

\vspace{0.6em}
\centering
\footnotesize
\renewcommand{\arraystretch}{1.3}
\resizebox{\linewidth}{!}{%
\begin{tabular}{r r r r r ccc}
& \textbf{price} & \textbf{yearOfReg.} & \textbf{powerPS} & \textbf{kilometer}
& \textbf{fuelType\_benzin} & \textbf{fuelType\_diesel} & \textbf{fuelType\_hybrid}\\
\hline
0 &  1450 & 1997 &  75 &  90000 & True  & False & False\\
1 & 13100 & 2005 & 280 &   5000 & True  & False & False\\
2 &  4500 & 2008 &  87 &  90000 & True  & False & False\\
3 &  6000 & 2009 & 177 & 125000 & False & True  & False\\
4 &  3990 & 1999 & 118 &  90000 & True  & False & False\\
\end{tabular}}

\source{\cite{seemann2024mlaz}, OneHotFahrzeugdaten.ipynb}
\end{frame}

\begin{frame}[fragile]{Ordinale Daten: Ersetzen}
\begin{columns}
\column{0.35\textwidth}
Bei ordinalen Daten lohnt sich zuerst ein \textbf{Überblick} über alle
vorkommenden Werte und ihre Häufigkeiten.

\vspace{0.5em}
\texttt{value\_counts()} liefert die Werte, \texttt{.count()} die Anzahl
verschiedener Kategorien.

\column{0.60\textwidth}
\begin{minted}[fontsize=\footnotesize]{python}
import pandas as pd
df = pd.read_csv("diamonds.csv")
df.head()

print(df["cut"].value_counts())
print(df["cut"].value_counts().count())
\end{minted}

\vspace{0.4em}
\begin{minted}[fontsize=\footnotesize, frame=none]{text}
cut
Ideal        21551
Premium      13791
Very Good    12082
Good          4906
Fair          1610
Name: count, dtype: int64
5
\end{minted}
\end{columns}

\source{\cite{seemann2024mlaz}, DiamantenOrdinaleDaten.ipynb}
\end{frame}

\begin{frame}[fragile]{Ordinale Daten: Ersetzen}
\footnotesize
\begin{columns}
\column{0.65\textwidth}
\begin{minted}[fontsize=\footnotesize]{python}
replace_map = {'cut': {'Fair': 1, 'Good': 2,
                       'Very Good': 3, 'Premium': 4,
                       'Ideal': 5}}
df = df.replace(replace_map)
df['cut'] = df['cut'].astype(int)
df.head()
\end{minted}

\column{0.3\textwidth}
Mit dem Wissen über die \textbf{Reihenfolge} erstellt man eine
\textbf{replace map} und übergibt sie an \texttt{replace()}.

\vspace{0.4em}
{\footnotesize\itshape Der Cast nach \texttt{int} ist nötig, da Pandas
nicht mehr automatisch umwandelt.}
\end{columns}

\vspace{0.6em}
\centering

\renewcommand{\arraystretch}{1.3}
\begin{tabular}{r r r c c c r r r r r}
& \textbf{carat} & \textbf{cut} & \textbf{color} & \textbf{clarity}
& \textbf{depth} & \textbf{table} & \textbf{price} & \textbf{x} & \textbf{y} & \textbf{z}\\
\hline
0 & 0.23 & 5 & E & SI2 & 61.5 & 55.0 & 326 & 3.95 & 3.98 & 2.43\\
1 & 0.21 & 4 & E & SI1 & 59.8 & 61.0 & 326 & 3.89 & 3.84 & 2.31\\
2 & 0.23 & 2 & E & VS1 & 56.9 & 65.0 & 327 & 4.05 & 4.07 & 2.31\\
3 & 0.29 & 4 & I & VS2 & 62.4 & 58.0 & 334 & 4.20 & 4.23 & 2.63\\
4 & 0.31 & 2 & J & SI2 & 63.3 & 58.0 & 335 & 4.34 & 4.35 & 2.75\\
\end{tabular}

\source{\cite{seemann2024mlaz}, DiamantenOrdinaleDaten.ipynb}
\end{frame}

\begin{frame}[fragile]{Unvollständige Daten: \texttt{dropna}}
\begin{columns}
\column{0.60\textwidth}
\begin{minted}[fontsize=\footnotesize]{python}
import pandas as pd
df = pd.read_csv("autos.csv")

# Anzahl fehlender Werte je Spalte
print(df.isna().sum())

# Zeilen mit fehlenden Werten entfernen
df = df.dropna()
df.head()
\end{minted}

\column{0.35\textwidth}
Reale Datensätze enthalten oft \textbf{fehlende Werte} (\texttt{NaN}).

\vspace{0.4em}
\texttt{dropna()} entfernt alle Zeilen mit mindestens einem fehlenden
Wert.

\vspace{0.4em}
{\footnotesize\itshape Alternativen: gezielt Spalten füllen
(\texttt{fillna}) statt löschen.}
\end{columns}

\vspace{0.6em}
\centering
{\footnotesize
\renewcommand{\arraystretch}{1.3}
\setlength{\tabcolsep}{5pt}
\begin{tabular}{r r r l r}
& \textbf{price} & \textbf{powerPS} & \textbf{fuelType} & \textbf{kilometer}\\
\hline
0 &  1450 &  75 & benzin & 90000\\
1 & 13100 & \cellcolor{red!18}NaN & benzin &  5000\\
2 &  4500 &  87 & \cellcolor{red!18}NaN & 90000\\
3 &  6000 & 177 & diesel & 125000\\
4 &  3990 & 118 & benzin & \cellcolor{red!18}NaN\\
\end{tabular}}

\vspace{0.4em}
{\footnotesize\itshape Zeilen 1, 2 und 4 werden durch \texttt{dropna()} entfernt.}
\end{frame}


\begin{frame}[fragile]{Unvollständige Daten: \texttt{fillna}}
\begin{columns}
\column{0.60\textwidth}
\begin{minted}[fontsize=\footnotesize]{python}
# Numerische Spalte mit Median fuellen
df["powerPS"] = df["powerPS"].fillna(
    df["powerPS"].median())

# Kategoriale Spalte mit häufigstem Wert
df["fuelType"] = df["fuelType"].fillna(
    df["fuelType"].mode()[0])
\end{minted}

\column{0.35\textwidth}
\texttt{fillna()} ersetzt fehlende Werte, statt Zeilen zu \textbf{löschen}.

\vspace{0.4em}
{\footnotesize\itshape Median für Zahlen (robust gegen Ausreißer), Modus
für Kategorien.}
\end{columns}

\vspace{0.5em}
\centering
{\footnotesize
\renewcommand{\arraystretch}{1.3}
\setlength{\tabcolsep}{6pt}
\begin{columns}
\column{0.48\textwidth}
\centering
\textbf{Vorher}\\[2pt]
\begin{tabular}{r r l}
& \textbf{powerPS} & \textbf{fuelType}\\
\hline
0 &  75 & benzin\\
1 & \cellcolor{red!18}NaN & benzin\\
2 &  87 & \cellcolor{red!18}NaN\\
3 & 177 & diesel\\
4 & 118 & benzin\\
\end{tabular}

\column{0.48\textwidth}
\centering
\textbf{Nachher}\\[2pt]
\begin{tabular}{r r l}
& \textbf{powerPS} & \textbf{fuelType}\\
\hline
0 &  75 & benzin\\
1 & \cellcolor{green!20}102 & benzin\\
2 &  87 & \cellcolor{green!20}benzin\\
3 & 177 & diesel\\
4 & 118 & benzin\\
\end{tabular}
\end{columns}}

\vspace{0.4em}
{\footnotesize\itshape Median(\texttt{powerPS})=102, Modus(\texttt{fuelType})=benzin.}

\end{frame}


\begin{frame}{ Übung Datenaufbereitung}
\begin{columns}
\column{0.60\textwidth}
Der Diamantendatensatz hat neben \texttt{cut} zwei weitere
\textbf{kategoriale} Spalten:

\begin{itemize}
  \item \textbf{color}: Werte D, E, F, \dots, J\\
        ({\itshape D am besten, J am schlechtesten})
  \item \textbf{clarity}: Reinheit, vom besten zum schlechtesten\\
        FL, IF, VVS1, VVS2, VS1, VS2, SI1, SI2, I1, I2, I3
\end{itemize}

\vspace{0.4em}
\textbf{Aufgabe:} Erweitern Sie Ihr Notebook zur Preisvorhersage, sodass
zusätzlich \textbf{color} und \textbf{clarity} berücksichtigt werden.
Bereiten Sie beide Spalten entsprechend auf.

\column{0.35\textwidth}
\centering
\twemoji[width=0.9\linewidth]{laptop}
\end{columns}
\source{Lösung: LösungDiamantenOrdinaleDatenKomplett.ipynb}
\end{frame}

\makequizpage{\QUIZURL}{\QUIZQRCODE}{\QUIZIMAGE}
\makequestionspage{\FRAGENIMAGE}